\begin{apartats}

\apartat{1}
Calcula els límits següents:
\begin{graella}{3}
  \sa \lim_{x\to3}\frac{x^2-9}{x^2-5x+6} &
  \sa \lim_{x\to-1}\frac{x^3+1}{x^2-1} &
  \sa \lim_{x\to2}\frac{x^2-4x+4}{x^2-3x+2}
\end{graella}

\begin{solucio}
Tots tres presenten una indeterminació $\tfrac00$: es factoritza i se simplifica.\\
i) $\dfrac{(x-3)(x+3)}{(x-3)(x-2)}=\dfrac{x+3}{x-2}\to\dfrac{6}{1}=6$.\\
ii) Per Ruffini, $x^3+1=(x+1)(x^2-x+1)$:
$\dfrac{(x+1)(x^2-x+1)}{(x+1)(x-1)}=\dfrac{x^2-x+1}{x-1}\to\dfrac{3}{-2}=-\dfrac32$.\\
iii) $\dfrac{(x-2)^2}{(x-2)(x-1)}=\dfrac{x-2}{x-1}\to\dfrac{0}{1}=0$.
\end{solucio}

\apartat{0,75}
Donada la funció $f(x)=\dfrac{2x}{\sqrt{x^2+5}}$, calcula:
\begin{graella}{3}
  \sa \lim_{x\to2}f(x) & \sa \lim_{x\to-1}f(x) & \sa \lim_{x\to0}f(x)
\end{graella}

\begin{solucio}
$f$ és contínua a tot $\mathbb{R}$, perquè el radicand és sempre positiu: n'hi ha prou
de substituir.\\
i) $\dfrac{4}{\sqrt9}=\dfrac43$. \quad
ii) $\dfrac{-2}{\sqrt6}=-\dfrac{2\sqrt6}{6}=-\dfrac{\sqrt6}{3}$. \quad
iii) $\dfrac{0}{\sqrt5}=0$.
\end{solucio}

\apartat{0,75}
Calcula, si existeixen, els límits següents. Justifica-ho amb els límits laterals.
\begin{graella}{2}
  \sa \lim_{x\to1}\frac{x+2}{(x-1)^2} & \sa \lim_{x\to1}\frac{x+2}{x-1}
\end{graella}

\begin{solucio}
En tots dos casos el numerador tendeix a $3$ i el denominador a $0$.\\
i) $(x-1)^2>0$ als dos costats de $1$: els dos laterals valen $+\infty$ i, per tant,
el límit és $+\infty$.\\
ii) $x-1$ canvia de signe: el lateral esquerre val $-\infty$ i el dret $+\infty$.
El límit \textbf{no existeix}.
\end{solucio}

\end{apartats}
